We have produced the greenhouse gas concentration forcing dataset for use in the historical, piControl
and other experiments in CMIP7 \citep[see publications that describe CMIP7 experiments, e.g.,][]{dunne_cmip7_2025}.
% ZNTODO: user guide link
A full user guide can be found at [TODO link].
Like the other forcings datasets, the latest version
and versions to use in CMIP experiments can be found at
% \newline \mbox{\url{https://input4mips-cvs.readthedocs.io/en/latest/}},
\url{https://input4mips-cvs.readthedocs.io/en/latest/},
with the specific page for greenhouse gas forcing being
\newline \mbox{\url{https://input4mips-cvs.readthedocs.io/en/latest/dataset-overviews/greenhouse-gas-concentrations/}}.
Further guidance on how to use forcings in each CMIP experiment can be found at
\newline \mbox{\url{https://wcrp-cmip.github.io/cmip7-guidance/docs/CMIP7/Experiment_set_up_and_Forcings/}}.

Compared to the CMIP6 dataset, the changes are minor.
% value-check: {"tag": "co2-diff-from-cmip6-all-radiative-effect", "unit": "W/m^2", "upper": 0.06, "lower": 0.055}
% value-check: {"tag": "co2-diff-from-cmip6-all-year", "unit": "yr", "upper": 250, "lower": 150}
% The other gases' differences are smaller, so CO_2's is the maximum
% value-check: {"tag": "ch4-diff-from-cmip6-all-radiative-effect", "unit": "W/m^2", "upper": 0.055, "lower": 0}
% value-check: {"tag": "n2o-diff-from-cmip6-all-radiative-effect", "unit": "W/m^2", "upper": 0.055, "lower": 0}
% value-check: {"tag": "cfc12-diff-from-cmip6-all-radiative-effect", "unit": "W/m^2", "upper": 0.055, "lower": 0}
% value-check: {"tag": "cfc12eq-diff-from-cmip6-all-radiative-effect", "unit": "W/m^2", "upper": 0.055, "lower": 0}
% value-check: {"tag": "hfc134aeq-diff-from-cmip6-all-radiative-effect", "unit": "W/m^2", "upper": 0.055, "lower": 0}
The maximum difference, in approximate radiative effect terms, is approximately 0.06~W/m2 in the CO_2 dataset around year 200.
Over the period used in the historical and piControl experiments, i.e. 1850 onwards,
% value-check: {"tag": "co2-diff-from-cmip6-1850-on-radiative-effect", "unit": "W/m^2", "upper": 0.03, "lower": 0.02}
% value-check: {"tag": "co2-diff-from-cmip6-1850-on-year", "unit": "yr", "upper": 1969, "lower": 1969}
% value-check: {"tag": "ch4-diff-from-cmip6-1850-on-radiative-effect", "unit": "W/m^2", "upper": 0.02, "lower": 0}
% value-check: {"tag": "n2o-diff-from-cmip6-1850-on-radiative-effect", "unit": "W/m^2", "upper": 0.02, "lower": 0}
% value-check: {"tag": "cfc12-diff-from-cmip6-1850-on-radiative-effect", "unit": "W/m^2", "upper": 0.02, "lower": 0}
% value-check: {"tag": "cfc12eq-diff-from-cmip6-1850-on-radiative-effect", "unit": "W/m^2", "upper": 0.02, "lower": 0}
% value-check: {"tag": "hfc134aeq-diff-from-cmip6-1850-on-radiative-effect", "unit": "W/m^2", "upper": 0.02, "lower": 0}
the maximum difference is approximately 0.025~W/m2 in the CO_2 dataset in 1969.

While the analysis we have done here suggests that the dataset
accurately and faithfully represents the variations in atmospheric greenhouse gas concentrations
over the last two thousand years, there are clear areas for improvement (Section~\ref{sec:discussion}).
To avoid delaying the wider CMIP work effort,
timely delivery was a key criteria for this dataset (Section~\ref{sec:output-dataset-requirements}),
which is why many elements were left unimproved
(and indeed why this manuscript comes out so long after the dataset itself was published).
We look foward to developing on this dataset to provide an improved product for users.
% % Include if we get confirmation of C3S before submission
% In particular, we are excited to provide annual updates of this and other forcing datasets
% as part of the European Union Horizon project Quality-assured Updates of forcing Information for Climate Change Assessments (QUICCA)
% and C3S project `'.
